% TOP_PROBLEM: {"id": 3600003, "title": "Multiplicity-one support of area Siegel–Veech constants", "category_id": 8, "category": {"id": 8, "name": "analysis", "display_name": "Analysis", "description": "Limits, continuity, calculus, and function theory.", "slug": "analysis", "order_index": 8, "created_at": "2026-07-31T15:26:25.671Z"}, "status": "open"}
% FIRST_POSTED: null
% ENTRY_KIND: statement_only
% Imported from: batch12.tex; UnsolvedMath ID 3600003
% Compile twice: lualatex 3600003.tex
\documentclass[11pt]{article}
\usepackage{fontspec}
\setmainfont{Latin Modern Roman}
\usepackage{amsmath,amssymb,mathtools,unicode-math}
\setmathfont{Latin Modern Math}
\usepackage{xurl,hyperref}
\hypersetup{hidelinks,pdftitle={UnsolvedMath Problem 3600003}}
\newcommand{\sref}[2]{\hyperlink{source:#1:#2}{[S#2]}}
\newcommand{\cataloguescope}[1]{\paragraph{Mathematical question.}#1}
\begin{document}
% UnsolvedMath ID 3600003; source catalogue code AMR-035-0003
\setcounter{section}{3600002}
\section{Multiplicity-one support of area Siegel–Veech constants}\label{3600003}
{\small\sffamily UnsolvedMath ID 3600003 · Analysis}
\subsection{Definitions and mathematical statement}

\paragraph{Shared notation.}$\mathbb N=\{1,2,\ldots\}$, and $\mathbb Z,\mathbb Q,\mathbb R,\mathbb C$ denote the integers, rationals, reals and complex numbers. Any other conventions are specified locally.


For an integer $g\ge2$, let $\widehat\Pi_{4g-4}$ consist of finite unordered lists $\boldsymbol d=(d_1,\ldots,d_n)$ with
\[
 d_i\in\{-1,0,1,2,\ldots\},\qquad\sum_i d_i=4g-4,\qquad
 \#\{i:d_i=-1\}\le\log g,
\]
where $\log$ is natural. The stratum $\mathcal Q_g(\boldsymbol d)$ parametrizes pairs $(X,q)$ consisting of a compact genus-$g$ Riemann surface and a meromorphic quadratic differential with the specified orders at its labeled singularities: $d_i=-1$ denotes a simple pole, $d_i>0$ a zero, and $d_i=0$ a marked regular point. Isomorphisms preserve the differential and labels. The quadratic stratum uses differentials that are not globally squares of Abelian differentials.

Fix the Masur--Veech normalization used by the sources. On the canonical orientation double cover $\widehat X$, where the pullback of $q$ is $\omega^2$, use anti-invariant relative periods of $\omega$. The period-coordinate lattice consists of linear forms taking values in $\mathbb Z+i\mathbb Z$ on anti-invariant integral relative homology, and its covolume is $1$. The relative marked set consists of preimages of zeros and marked regular points. Put $\mathcal Q_{g,1}(\boldsymbol d)=\{(X,q):\int_X|q|=1/2\}$. If $D=2g+n-2$ is the complex dimension, define $\mu_1$ by
\[
 d\mu=r^{2D-1}\,dr\,d\mu_1,\qquad q=r^2q_1,\quad q_1\in\mathcal Q_{g,1}(\boldsymbol d),
\]
and $\operatorname{Vol}\mathcal Q_g(\boldsymbol d)=\int_{\mathcal Q_{g,1}(\boldsymbol d)}d\mu_1$. In particular zeros are labeled; adding a labeled regular point multiplies the volume by $2$. Only nonempty strata enter the questions.
For a connected component $\mathcal C$ of the unit-area slice, its area Siegel--Veech constant $c_{\mathrm{area}}(\mathcal C)$ counts maximal flat cylinders with weight equal to cylinder area divided by total surface area. More explicitly, for almost every surface $S\in\mathcal C$, of area $A$, put
\[
N_{\mathrm{area}}(S,L)=\sum_{\substack{Z\text{ maximal cylinder on }S\\\text{circumference}(Z)\le L}}
\frac{\operatorname{Area}(Z)}{A},\qquad
N_{\mathrm{area}}(S,L)\sim c_{\mathrm{area}}(\mathcal C)\,\frac{\pi L^2}{A}.
\]
Almost every refers to the normalized Masur--Veech measure on the component. A non-hyperelliptic component is a connected component other than a hyperelliptic component in the stratum classification.
A saddle connection is a flat geodesic joining singularities (possibly the same one), with no singularity in its interior. Two such connections are $\widehat{\text{homologous}}$ when their anti-invariant lift classes on the orientation double cover agree up to orientation, using the convention of Goujard §2.1: for lifts $\gamma',\gamma''$, use $[\widehat\gamma]=[\gamma']$ if $[\gamma']+[\gamma'']=0$, and otherwise use $[\widehat\gamma]=[\gamma']-[\gamma'']$. A configuration records a maximal such collection, the genera and singularities of the complementary pieces, their angular sectors, and the gluing.

The two selected multiplicity-one cylinder configurations are those depicted in Conjecture 3 of the source, described as follows. In both, a single cylinder has one saddle-connection loop on each boundary; its non-cylinder complement is connected and has nontrivial linear holonomy.
\begin{itemize}
\item $\mathcal C_{b,\mathrm I}(d_i,d_j)$ has the two boundary loops based at distinct labeled zeros of orders $d_i,d_j\ge1$. Shrinking the cylinder produces the genus-$(g-1)$ principal boundary stratum with these orders replaced by $d_i-2,d_j-2$.
\item $\mathcal C_{b,\mathrm{II}}(a_1,a_2)$ has both boundary loops based at the same labeled zero of order $a_1+a_2+2$. The two complementary angular sectors there have angles $(a_1+1)\pi$ and $(a_2+1)\pi$. Its principal boundary stratum has genus $g-1$ and replaces this order by $a_1+a_2-2$.
\end{itemize}
Require $a_1,a_2\in\{0,1,\ldots\}$ and $a_1+a_2\ge3$, exactly as in the conjecture. Let $\mathfrak B(\boldsymbol d)$ be the set of all distinct admissible configurations of these two types, with labeled zeros as specified; reversing a direction or exchanging the two descriptions of the same geometric cylinder does not count it twice. The contribution $c_{\mathrm{area}}(\mathcal B)$ is defined by the same weighted counting limit as above, restricted to cylinders realizing $\mathcal B$. If the stratum is disconnected, use the Masur--Veech-volume-weighted average of its component constants.
\paragraph{Statement recorded in the supplied source.}
Prove that the listed configurations asymptotically account for the whole area constant: along every sequence of nonempty strata $\mathcal Q_g(\boldsymbol d)$ with $\boldsymbol d\in\widehat\Pi_{4g-4}$ and $g\to\infty$,
\[
 \frac{\displaystyle\sum_{\mathcal B\in\mathfrak B(\boldsymbol d)}c_{\mathrm{area}}(\mathcal B)}
 {c_{\mathrm{area}}(\mathcal Q_g(\boldsymbol d))}\longrightarrow1.
\]
This is the full-stratum support question, with no non-hyperelliptic restriction added. \sref{3600003}{2} \sref{3600003}{3} \sref{3600003}{4}
\subsection{Short English statement}
Show that, in the specified large-genus quadratic strata, almost all of the area-weighted cylinder contribution comes from the two stated single-cylinder configurations.
\subsection{Sources}
\begin{description}
\item[\hypertarget{source:3600003:1}{S1}] ULAM.AI and the UnsolvedMath Contributors, UnsolvedMath: A Curated Collection of Open Mathematics Problems, record 3600003 (AMR-035-0003). CC BY 4.0. \url{https://huggingface.co/datasets/ulamai/UnsolvedMath}.
\item[\hypertarget{source:3600003:2}{S2}] Amol Aggarwal, Vincent Delecroix, Élise Goujard, Peter Zograf and Anton Zorich, Conjectural Large Genus Asymptotics of Masur–Veech Volumes and of Area Siegel–Veech Constants of Strata of Quadratic Differentials (2019), Conjecture 3 and its diagram; §3, pp.7–9. \url{https://arxiv.org/abs/1912.11702}.
\item[\hypertarget{source:3600003:3}{S3}] Jayadev S. Athreya, Alex Eskin and Anton Zorich, Right-angled billiards and volumes of moduli spaces of quadratic differentials on CP1 (2016), §4.1, Convention 4.1, equations (4.2)–(4.5); §4.13. \url{https://arxiv.org/abs/1212.1660}.
\item[\hypertarget{source:3600003:4}{S4}] Élise Goujard, Siegel–Veech constants for strata of moduli spaces of quadratic differentials (2015), §§1.1–1.3,2.1–2.3,3.1; Figure 10. \url{https://arxiv.org/abs/1405.5899}.
\end{description}
\label{end:3600003}
\end{document}
